% © 2026 Ing. David Jaroš — CC BY-NC-ND 4.0
%
% File: research_tracks/T3_ALPHA/information_loss_alpha_self_consistency.tex
% Status: RESEARCH TRACK / HYPOTHESIS. Not canonical. Not a proof of alpha.
% Purpose: Spectral information-loss correction to the eta-winding alpha route,
%          with Layer2 observable-sector projection interpretation.
% Date: 2026-06-12

\documentclass[12pt,a4paper]{article}
\usepackage[a4paper,margin=2.5cm]{geometry}
\usepackage{amsmath,amssymb,amsthm,mathtools}
\usepackage{booktabs}
\usepackage{hyperref}
\usepackage{xcolor}
\usepackage[most]{tcolorbox}
\usepackage{microtype}

\hypersetup{
  colorlinks=true,
  linkcolor=blue!60!black,
  citecolor=green!50!black,
  urlcolor=blue!60!black
}

\newtheorem{definition}{Definition}[section]
\newtheorem{proposition}[definition]{Proposition}
\newtheorem{conjecture}[definition]{Conjecture}
\newtheorem{remark}[definition]{Remark}

\newtcolorbox{statusbox}[1][]{
  colback=gray!8!white,
  colframe=black!60,
  arc=3pt,
  boxrule=1pt,
  title=Status Summary,
  #1
}

\title{A Layer2 Projection Correction to the UBT Alpha Information-Loss Route}
\author{Ing.\ David Jaro\v{s}}
\date{2026-06-12}

\begin{document}
\maketitle

\begin{abstract}
This research note refines the UBT fine-structure constant route by combining
the Layer1 eta/theta spectrum with the Layer2 observable-sector
coding/projection layer.  The eta-winding coefficient is
\[
  B(\rho)=12^{3/2}\bigl(2\eta(i\rho)\bigr)^{1/4},
\]
where $\rho$ is the effective quasi-periodic theta-modulus ratio.  The
stationary winding index $n$ is defined by
\[
  \frac{2n}{\ln n+1}=B(\rho).
\]
The uncorrected self-dual value $\rho=1$ gives
$n=136.989099634\ldots$.  A minimal four-channel information-loss ansatz gives
$n=137.0368227057\ldots$.  The present note introduces a systematic
eta-spectral Layer2 projection correction
\[
  C_Q(n)
  =
  4-\frac{\pi-3}{2}\frac{n-3}{n}.
\]
Here
\[
  \frac{\pi-3}{2}
  =
  12\pi
  \sum_{m=1}^{\infty}\frac{m}{e^{2\pi m}-1}
\]
is the self-dual eta-mode occupation subtraction, while $(n-3)/n$ is
interpreted as the finite-rank complement of a protected three-dimensional
Layer2 SU(3)/color-code subspace.  The resulting closed equation gives
\[
  n_{\rm UBT}=137.035999142931\ldots .
\]
Relative to CODATA/NIST 2022,
$\alpha^{-1}=137.035999177(21)$, this is lower by approximately $1.62\sigma$.
The result remains research-track: the remaining theorem is to derive the
Layer2 projection-rank factor from the UBT projection kernel $\Pi_{\rm L2}$.
\end{abstract}

\begin{statusbox}
\textbf{Status.} Research-track only.  This note does not claim a canonical
derivation of the fine-structure constant.  It replaces a loose entropy tuning
by a Layer1/Layer2 mechanism: eta spectral subtraction from Layer1 and
finite-rank observable-sector projection from Layer2.  The explicit remaining
gap is the derivation of the projection-rank factor $(n-3)/n$ from the UBT
Layer2 projector.
\end{statusbox}

\section{Separation of indices}

We distinguish three quantities:
\[
  m \in \mathbb{N}
\]
for eta/theta product modes,
\[
  n=\alpha^{-1}
\]
for the effective stationary winding index, and
\[
  \rho=\frac{R_\tau}{R_\psi}
\]
for the effective quasi-periodic theta-modulus ratio.

The eta/theta product index $m$ is not identified with the physical value
$n\approx137$.  Instead, the full eta spectrum contributes to the coefficient
$B(\rho)$, and the stationary equation then determines the effective physical
index $n$.

\section{Layer1 eta spectrum}

For $\tau=i\rho$ the Dedekind eta function is
\[
  \eta(i\rho)
  =
  e^{-\pi\rho/12}
  \prod_{m=1}^{\infty}
  \left(1-e^{-2\pi m\rho}\right).
\]
The logarithm is
\[
  L(\rho)=\log\eta(i\rho)
  =
  -\frac{\pi\rho}{12}
  +
  \sum_{m=1}^{\infty}
  \log\left(1-e^{-2\pi m\rho}\right).
\]
Differentiation gives
\[
  \frac{dL}{d\rho}
  =
  -\frac{\pi}{12}
  +
  2\pi
  \sum_{m=1}^{\infty}
  \frac{m}{e^{2\pi m\rho}-1}.
\]
The denominator is $e^{2\pi m\rho}-1$.

At the self-dual point $\rho=1$ the Eisenstein identity
\[
  E_2(i)=\frac{3}{\pi}
\]
implies
\[
  \frac{d}{d\rho}\log\eta(i\rho)\bigg|_{\rho=1}
  =
  -\frac14.
\]
Equivalently,
\[
  \sum_{m=1}^{\infty}
  \frac{m}{e^{2\pi m}-1}
  =
  \frac{1}{24}-\frac{1}{8\pi}.
\]
Therefore
\[
  12\pi
  \sum_{m=1}^{\infty}
  \frac{m}{e^{2\pi m}-1}
  =
  \frac{\pi-3}{2}.
\]

\begin{remark}
The factor $(\pi-3)/2$ is not a fitted decimal.  It is the self-dual eta-mode
occupation subtraction derived from the Layer1 eta spectrum.
\end{remark}

\section{Eta-winding stationary equation}

The eta-corrected winding coefficient is
\[
  B(\rho)
  =
  12^{3/2}
  \left(2\eta(i\rho)\right)^{1/4}.
\]
The effective winding potential is modeled so that its stationary condition is
\[
  \frac{2n}{\ln n+1}=B(\rho).
\]
At the uncorrected self-dual point $\rho=1$, this gives
\[
  n_0=136.989099634108\ldots .
\]

\section{Information loss and the four-channel baseline}

The effective quasi-periodic modulus is written as
\[
  \rho=\exp(-\Delta I_Q),
\]
where $\Delta I_Q$ is the average information loss per quasi-periodic theta
cycle.

The coarse one-loop ansatz is
\[
  \Delta I_Q(n)=\frac{C_Q}{2\pi n}.
\]
The minimal Dirac-biquaternionic channel count suggests
\[
  C_Q=4.
\]
The resulting closed equation is
\[
  \frac{2n}{\ln n+1}
  =
  12^{3/2}
  \left[
    2\eta\!\left(
      i\exp\!\left[-\frac{4}{2\pi n}\right]
    \right)
  \right]^{1/4}.
\]
Numerically,
\[
  n=137.036822705722871\ldots .
\]
This value is close but too high.

\section{Layer2 observable-sector projection}

Layer2 is the UBT observable-sector coding/projection layer.  It should not be
introduced as an external ``simulation assumption'' in the derivation.
Mathematically it is a finite-rank projection layer selecting physical
representation subspaces from the larger theta-biquaternionic field.

In the QCD/SU(3) sector the Layer2 construction uses a three-dimensional
protected color-code subspace.  Abstractly we write
\[
  \mathcal H_{\rm color}
  =
  \mathrm{span}\{|100\rangle,\ |010\rangle,\ |001\rangle\},
\]
with projector
\[
  \Pi_{\rm color},
  \qquad
  \mathrm{rank}\,\Pi_{\rm color}=3.
\]
In a finite effective winding sector of dimension $n$, the eta-spectral
information-loss operator acts on the complement of this protected subspace:
\[
  I_n-\Pi_{\rm color}.
\]
The normalized rank of this complement is
\[
  \frac{\mathrm{Tr}(I_n-\Pi_{\rm color})}{\mathrm{Tr}(I_n)}
  =
  \frac{n-3}{n}.
\]

\begin{conjecture}[Layer2 projection-rank correction]
The eta-spectral information-loss subtraction acts only on the complement of
the protected Layer2 color-code subspace.  Therefore the effective subtraction
of the four-channel coefficient is
\[
  4-C_Q(n)
  =
  \varepsilon_\eta\frac{n-3}{n},
\]
where
\[
  \varepsilon_\eta
  =
  12\pi
  \sum_{m=1}^{\infty}
  \frac{m}{e^{2\pi m}-1}
  =
  \frac{\pi-3}{2}.
\]
Thus
\[
  C_Q(n)
  =
  4-\frac{\pi-3}{2}\frac{n-3}{n}.
\]
\end{conjecture}

\begin{remark}
The factor $(n-3)/n$ is not claimed as a theorem in this note.  Its proposed
origin is the Layer2 projection rank: three protected SU(3)/color-code states
are removed from the eta-loss sector.  The remaining proof gap is to derive
this projector from the canonical UBT Layer2 construction.
\end{remark}

\section{Closed Layer1/Layer2 self-consistency equation}

Combining the eta-winding equation with the Layer2 projection correction gives
\[
  \boxed{
  \frac{2n}{\ln n+1}
  =
  12^{3/2}
  \left[
    2\eta\!\left(
      i\exp\!\left[
        -\frac{1}{2\pi n}
        \left(
          4-\frac{\pi-3}{2}\frac{n-3}{n}
        \right)
      \right]
    \right)
  \right]^{1/4}.
  }
\]
Solving gives
\[
  \boxed{
  n_{\rm UBT}=137.035999142931005\ldots .
  }
\]

For comparison:
\[
\begin{array}{lll}
\toprule
\text{Model} & \text{Input} & n \\
\midrule
\text{Self-dual eta-winding} & \rho=1 & 136.989099634108\ldots \\
\text{Four-channel information loss} & C_Q=4 & 137.036822705723\ldots \\
\text{Layer2 eta-spectral projection} &
C_Q(n)=4-\frac{\pi-3}{2}\frac{n-3}{n}
& 137.035999142931\ldots \\
\bottomrule
\end{array}
\]

Using CODATA/NIST 2022
\[
  \alpha^{-1}_{\rm CODATA}=137.035999177(21),
\]
the Layer2 eta-spectral projection value differs by
\[
  -3.4069\times10^{-8},
\]
or approximately
\[
  -1.62\sigma.
\]

\section{Interpretation}

The result can be summarized as a Layer1/Layer2 fixed point:
\[
  \text{Layer1 eta/theta spectrum}
  \quad+\quad
  \text{Layer2 observable-sector projection}
  \quad\Rightarrow\quad
  n=\alpha^{-1}.
\]
The modular eta spectrum supplies the universal subtraction
\[
  \varepsilon_\eta=\frac{\pi-3}{2},
\]
and the Layer2 color-code projection supplies the finite-rank factor
\[
  \frac{n-3}{n}.
\]
Thus the fine-structure constant is not identified with a single theta mode.
Rather, it appears as the stationary effective winding index generated by the
full eta spectrum after Layer2 observable-sector projection.

\section{Remaining proof gap}

The remaining open step is structural:
\[
  \boxed{
  G137\text{-L2}:
  \quad
  \text{derive }
  \frac{\mathrm{Tr}(I_n-\Pi_{\rm color})}{\mathrm{Tr}I_n}
  =
  \frac{n-3}{n}
  \text{ from canonical Layer2 UBT.}
  }
\]
This requires explicitly constructing the Layer2 projector
\[
  \Pi_{\rm L2}:
  \mathbb B\otimes\mathcal H_\Theta
  \longrightarrow
  \mathcal H_{\rm obs}
\]
and showing that the three-dimensional SU(3)/color-code subspace is protected
from eta-spectral information loss.

Until this is done, the result is a systematic research-track derivation with
one explicit Layer2 projection theorem remaining, not a canonical derivation
of $\alpha$.

\section{Conclusion}

This note replaces the earlier heuristic entropy refinement by a more
structured Layer1/Layer2 mechanism.  The final research-track equation is
\[
  \frac{2n}{\ln n+1}
  =
  12^{3/2}
  \left[
    2\eta\!\left(
      i\exp\!\left[
        -\frac{1}{2\pi n}
        \left(
          4-\frac{\pi-3}{2}\frac{n-3}{n}
        \right)
      \right]
    \right)
  \right]^{1/4}.
\]
It predicts
\[
  \alpha^{-1}_{\rm UBT}=137.035999142931\ldots ,
\]
approximately $1.62\sigma$ below the CODATA/NIST 2022 value.  The result is
strong enough to justify further work, but the Layer2 projection-rank theorem
must be proved before promoting the alpha route to the canonical layer.

\end{document}
